% ============================================================================
%  Préambule — mise en forme, paquets, macros
% ============================================================================

% ---- Langue et encodage -----------------------------------------------------
\usepackage[T1]{fontenc}
\usepackage[utf8]{inputenc}
\usepackage{lmodern}
\usepackage[french]{babel}
\usepackage{microtype}

% ---- Mise en page -----------------------------------------------------------
\usepackage[a4paper,top=2.8cm,bottom=2.8cm,left=2.7cm,right=2.7cm]{geometry}
\usepackage{setspace}
\onehalfspacing
\usepackage{emptypage}

% ---- Mathématiques ----------------------------------------------------------
\usepackage{amsmath,amssymb,amsthm}
\usepackage{mathtools}
\usepackage{bm}
\usepackage[version=4]{mhchem}   % formules chimiques : \ce{H3PO4}

% ---- Tableaux et figures ----------------------------------------------------
\usepackage{booktabs}
\usepackage{array}
\usepackage{multirow}
\usepackage{longtable}
\usepackage{graphicx}
\usepackage{float}
\usepackage{caption}
\captionsetup{font=small,labelfont=bf}

% ---- Nombres ----------------------------------------------------------------
\usepackage{siunitx}
\sisetup{
  output-decimal-marker = {,},
  group-separator = {\,},
  group-minimum-digits = 4,
  detect-weight = true,
}

% ---- Couleurs et encadrés ---------------------------------------------------
\usepackage{xcolor}
\definecolor{ocpvert}{RGB}{0,124,84}
\definecolor{ocpbleu}{RGB}{31,78,121}
\definecolor{grisclair}{RGB}{244,246,248}
\definecolor{grisfonce}{RGB}{90,95,100}
\definecolor{rougealerte}{RGB}{176,42,42}

\usepackage[most]{tcolorbox}

% ---- Schémas ----------------------------------------------------------------
\usepackage{tikz}
\usetikzlibrary{shapes.geometric,arrows.meta,positioning,calc,fit,backgrounds}

% ---- Code -------------------------------------------------------------------
\usepackage{listings}
\lstset{
  basicstyle=\ttfamily\footnotesize,
  keywordstyle=\color{ocpbleu}\bfseries,
  commentstyle=\color{grisfonce}\itshape,
  stringstyle=\color{ocpvert},
  numbers=left,
  numberstyle=\tiny\color{grisfonce},
  numbersep=8pt,
  frame=leftline,
  framesep=6pt,
  rulecolor=\color{ocpbleu},
  breaklines=true,
  showstringspaces=false,
  columns=fullflexible,
  extendedchars=true,
  literate=%
    {é}{{\'e}}1 {è}{{\`e}}1 {ê}{{\^e}}1 {à}{{\`a}}1 {ç}{{\c c}}1
    {É}{{\'E}}1 {î}{{\^i}}1 {ô}{{\^o}}1 {û}{{\^u}}1 {ù}{{\`u}}1,
}

% ---- Titres -----------------------------------------------------------------
\usepackage{titlesec}
\titleformat{\chapter}[display]
  {\normalfont\huge\bfseries\color{ocpbleu}}
  {\normalsize\normalfont\color{grisfonce}\MakeUppercase{\chaptertitlename}\ \thechapter}
  {14pt}{\Huge}
\titlespacing*{\chapter}{0pt}{-20pt}{30pt}
\titleformat{\section}{\normalfont\Large\bfseries\color{ocpbleu}}{\thesection}{1em}{}
\titleformat{\subsection}{\normalfont\large\bfseries\color{grisfonce}}{\thesubsection}{1em}{}

% ---- En-têtes et pieds ------------------------------------------------------
\usepackage{fancyhdr}
\pagestyle{fancy}
\fancyhf{}
\fancyhead[L]{\small\color{grisfonce}\nouppercase{\leftmark}}
\fancyhead[R]{\small\color{grisfonce}\thepage}
\renewcommand{\headrulewidth}{0.4pt}
\fancypagestyle{plain}{\fancyhf{}\fancyhead[R]{\small\color{grisfonce}\thepage}
  \renewcommand{\headrulewidth}{0pt}}

% ---- Liens ------------------------------------------------------------------
\usepackage[hidelinks,breaklinks]{hyperref}
\hypersetup{
  colorlinks = true,
  linkcolor  = ocpbleu,
  citecolor  = ocpvert,
  urlcolor   = ocpbleu,
  pdftitle   = {Optimisation de la production et de la distribution d'acide phosphorique},
  pdfsubject = {Rapport de stage d'ingénieur},
}
\usepackage[french]{cleveref}

% ============================================================================
%  Environnements personnalisés
% ============================================================================

\newtcolorbox{encadre}[1][]{
  enhanced, breakable,
  colback=grisclair, colframe=ocpbleu,
  boxrule=0pt, leftrule=3pt,
  arc=0pt, outer arc=0pt,
  left=10pt, right=10pt, top=8pt, bottom=8pt,
  fonttitle=\bfseries, #1
}

\newtcolorbox{resultat}[1][]{
  enhanced, breakable,
  colback=ocpvert!5, colframe=ocpvert,
  boxrule=0pt, leftrule=3pt,
  arc=0pt, outer arc=0pt,
  left=10pt, right=10pt, top=8pt, bottom=8pt,
  fonttitle=\bfseries\color{ocpvert}, #1
}

\newtcolorbox{alerte}[1][]{
  enhanced, breakable,
  colback=rougealerte!5, colframe=rougealerte,
  boxrule=0pt, leftrule=3pt,
  arc=0pt, outer arc=0pt,
  left=10pt, right=10pt, top=8pt, bottom=8pt,
  fonttitle=\bfseries\color{rougealerte}, #1
}

\theoremstyle{definition}
\newtheorem{definitionn}{Définition}[chapter]
\theoremstyle{plain}
\newtheorem{proposition}{Proposition}[chapter]

% ============================================================================
%  Macros de notation du modèle
% ============================================================================

\newcommand{\Lvingtneuf}{\ensuremath{\mathcal{L}_{29}}}
\newcommand{\Lcinqquatre}{\ensuremath{\mathcal{L}_{54}}}
\newcommand{\Lvingtneufetoile}{\ensuremath{\mathcal{L}^{\ast}_{29}}}
\newcommand{\Ech}{\ensuremath{\mathcal{E}}}
\newcommand{\Conso}{\ensuremath{\mathcal{K}}}
\newcommand{\Acides}{\ensuremath{\mathcal{A}}}
\newcommand{\Prod}{\ensuremath{\mathcal{P}}}

\newcommand{\xstd}{\ensuremath{x^{\mathrm{std}}}}
\newcommand{\xdec}{\ensuremath{x^{\mathrm{dec}}}}
\newcommand{\uncl}{\ensuremath{u^{\mathrm{ncl}}}}
\newcommand{\udec}{\ensuremath{u^{\mathrm{dec}}}}
\newcommand{\etacl}{\ensuremath{\eta^{\mathrm{cl}}}}
\newcommand{\etacoc}{\ensuremath{\eta^{\mathrm{coc}}}}
\newcommand{\etaconc}{\ensuremath{\eta^{\mathrm{conc}}}}
\newcommand{\Picoc}{\ensuremath{\Pi^{\mathrm{coc}}}}
\newcommand{\Pincl}{\ensuremath{\Pi^{\mathrm{ncl}}}}

\newcommand{\PdeuxOcinq}{\ensuremath{\mathrm{P_2O_5}}}
\newcommand{\tP}{\,t~\PdeuxOcinq}

% Renvoi vers une contrainte du modèle (nom court : \Cref appartient à cleveref)
\newcommand{\Ctr}[1]{\textbf{(C#1)}}
